\addchap*{Abstract}
Large Language Models (LLMs) have been increasingly studied for knowledge-intensive tasks across different domains such as finance, medicine, law, scientific discovery and open-domain question answering. Retrieval-augmented generation (RAG) techniques can improve generation quality and reduce hallucinations by providing external knowledge to LLMs. However, both conventional RAG and general LLM querying may be less effective for complex domains containing narrative-rich information, as well as for database-style queries such as SELECT, COUNT, and ORDER BY. The main goal of this thesis is to investigate whether providing LLMs with structured data containing narrative information can improve answer quality and reduce hallucinations in the domain of historical battles. First, we construct a large-scale knowledge base of 11,561 historical battles identified via Wikidata. Each battle is represented in a structured JSON file containing metadata, military facets, narrative information and patterns, generated from the raw text of the corresponding Wikipedia article using the GPT-OSS-120B model. The structured battle representations are then retrieved using three different retrieval approaches and provided as context to the same model. During the experimental part, we compare the question-answering performance of the baseline LLM and context-augmented LLM using several evaluation metrics. The main contribution of this thesis is to design a structured representation and retrieval framework for the historical battles domain and analyze how structured context affects LLM-based question-answering performance compared with a baseline LLM.